% =====================================================================
% Appendix: supplementary figures
%
% Collects the figures that are produced by the analysis code but are not
% reproduced in the main chapters. It is designed to grow one section per
% chapter; for now it contains only the dataset figures.
%
% Placement: after the last chapter of the main matter. \appendix must appear
% exactly once in the document -- if it is already declared elsewhere, delete
% the \appendix line below. If the thesis class has no \chapter (e.g. article),
% demote \chapter -> \section and each \section -> \subsection.
%
% Figures: docs/images/dataset/ in the repository, expected at images/datasets/
% in the thesis tree. All are exported by analysis/datasets_analysis.ipynb at
% their final printed width and must be included with width=\linewidth without
% further scaling.
%
% Editorial rule applied here: every figure keeps a self-contained caption that
% states what is plotted and how to read the encoding, plus one short paragraph
% giving the fact it supports and the place in the main text where that fact is
% used. A figure that cannot be read without the main text does not belong in an
% appendix, and an appendix of bare plots cannot be assessed.
% =====================================================================

\appendix

\chapter{Supplementary figures}
\label{app:figures}

This appendix collects the figures produced during the analysis that are not reproduced in the main chapters. They are included for completeness and reproducibility: each was generated by the same code and from the same data as the corresponding main-text material, and each is referenced from the chapter it belongs to. A figure appears here rather than in the main text when it is redundant with a table or with another figure, when it supports a statement that a single number already conveys, or when it is qualitative and therefore cannot carry an argument on its own. Nothing in this appendix is required to follow the main line of reasoning.


\section{Datasets}
\label{app:figures_datasets}

The four figures below complement Section~\ref{sec:dataset_analysis}. All are produced by the dataset analysis notebook from the same adapters used at training time, so split sizes, class indices, and hierarchy edges are identical to those seen by the models. As in the main text, they are descriptive: they characterize the data and the label space, not the behaviour of any model.

\subsection{Split sizes and composition}

Figure~\ref{fig:app_split_composition} shows the same partition reported numerically in Table~\ref{tab:protocol_splits}. It is kept out of the main text because nine numbers are read more precisely from the table than from a bar chart, but the percentage panel makes one asymmetry immediate: CIFAR-100 devotes 75\% of its images to training and only 8\% to validation, CUB-200-2011 44\% and 7\%, whereas the official FGVC-Aircraft partition splits the dataset into equal thirds. Model selection is therefore performed on very differently sized validation sets across the three benchmarks.

\begin{figure}[ht]
\centering
\includegraphics[width=\linewidth]{images/datasets/dataset_scale_and_taxonomy_width.pdf}
\caption{Partition used by the unified protocol. Panel (a) gives absolute sizes: each bar is one dataset and its segments are the mutually exclusive train, validation, and test splits, so the bar height is the total number of images. Panel (b) repeats the same three splits as percentages of each dataset, which keeps the composition of the two smaller benchmarks readable next to CIFAR-100.}
\label{fig:app_split_composition}
\end{figure}

\subsection{Distribution of the branching factor}

Table~\ref{tab:branching} reports the median, maximum, and number of single-child parents of each transition; Figure~\ref{fig:app_branching} shows the full distribution behind those summaries. It confirms that the summary statistics are not hiding a different shape: the CUB family-to-species transition is genuinely long-tailed, with most families holding one to four species and a single family holding thirty, while the FGVC-Aircraft family-to-variant transition is concentrated almost entirely on one child per parent. The CIFAR-100 lower transition appears as a single horizontal line at five, since every superclass contains exactly five fine classes.

\begin{figure}[ht]
\centering
\includegraphics[width=\linewidth]{images/datasets/hierarchy_branching_beeswarms.pdf}
\caption{Number of immediate children per parent class, for both transitions of each taxonomy. Every circle is one parent class and its vertical position is the exact integer number of children; parents sharing a value are spread symmetrically around the transition centre, so horizontal width encodes how many parents share that fan-out. The black segment is the median and the white diamond the mean. The annotation reports the number of parents with a single child, for which the level pair carries no decision.}
\label{fig:app_branching}
\end{figure}

\subsection{Leaf classes covered by each top-level class}

Figure~\ref{fig:app_descendants} quantifies per class what the top band of the icicle plot (Figure~\ref{fig:taxonomy_icicle}) shows as areas. It is the more precise of the two views and the less immediate one, which is why the icicle is used in the main text. It makes the imbalance of the CUB taxonomy explicit: the largest order covers 133 of the 200 species and the second 23, after which every remaining order covers seven or fewer. CIFAR-100 decays gently from 25 to 5 fine classes per coarse-1 group, and FGVC-Aircraft from 22 to 1 variant per manufacturer, with 13 of its 30 manufacturers covering a single variant.

\begin{figure}[ht]
\centering
\includegraphics[width=\linewidth]{images/datasets/leaf_descendants_per_root.pdf}
\caption{Number of leaf classes reachable beneath each top-level class, following both transitions. Bars are sorted by size, so the horizontal position is a rank and not a class identifier, and positions must not be matched across datasets. The annotation names the largest top-level class and its share of the leaf classes.}
\label{fig:app_descendants}
\end{figure}

\subsection{Representative images}

Figure~\ref{fig:app_gallery} is qualitative and is included only to make the difference in visual regime concrete. It illustrates why the three benchmarks are not interchangeable: CIFAR-100 classes are recognizable from coarse colour and shape at $32\times32$, whereas distinguishing two CUB species or two Aircraft variants depends on localized detail --- plumage patterning, engine count and placement, tail geometry --- that survives only at higher resolution. It does not measure visual diversity, inter-class similarity, background bias, or intrinsic difficulty, and no claim in this work rests on it.

\begin{figure}[ht]
\centering
\includegraphics[width=\linewidth]{images/datasets/representative_dataset_images.png}
\caption{One training image from each of six leaf classes per dataset, drawn with a fixed seed. Each row is one dataset and each column an independently selected class, so column positions do not correspond to related concepts across rows. Titles give the leaf-class name. Images are shown at their native aspect ratio, not at the resolution seen by the models.}
\label{fig:app_gallery}
\end{figure}
